\section*{Capítulo \ref{ch:b3:quantitative-geomorphology} --- Geomorfología cuantitativa}

\begin{solution}{exo:b3:quantitative-geomorphology:1}
$a/\sqrt{\pi\kappa t} = 0.30 - 0.04 = 0.26$, por lo que $\kappa t = (2/0.26)^2/\pi =
\qty{18.8}{m^2}$, y $t = 18.8/(5\times10^{-3}) = \qty{3.8}{kyr}$, dentro de un
factor de dos en cualquier sentido, dado el rango de $\kappa$.
\end{solution}

\begin{solution}{exo:b3:quantitative-geomorphology:2}
$E = \kappa\,|\nabla^2 z| = 4\times10^{-3}\times0.012 =
\qty{4.8e-5}{m/yr}$, \qty{0.048}{mm/yr}. El relieve $E L^2/(2\kappa) =
0.012\times80^2/2 = \qty{38.4}{m}$.
\end{solution}

\begin{solution}{exo:b3:quantitative-geomorphology:3}
$E = 2\times10^{-6}\times(25\times10^6)^{0.5}\times0.01 = 2\times10^{-6}
\times5000\times0.01 = \qty{1.0e-4}{m/yr}$, \qty{0.1}{mm/yr}. Como $E$ está
en \unit{m/yr}, $A^m$ en $\unit{m}^{2m}$ y $S$ no tiene unidades, $K$ está en
$\unit{m}^{1-2m}\,\unit{yr^{-1}}$: \unit{yr^{-1}} para $m = 0.5$.
\end{solution}

\begin{solution}{exo:b3:quantitative-geomorphology:4}
(a) Fija: un lecho duro empina el río donde aflora, y el escalón
permanece con el lecho. (b) Anclada en la falla: mientras la falla siga
desplazándose, el río se empina al cruzarla. (c) Móvil: el levantamiento de
la costa rebajó a la vez el \omterm{def:b3:quantitative-geomorphology:knickpoint}{nivel de base} de todos los ríos; los \omterm{def:b3:quantitative-geomorphology:knickpoint}{puntos de
ruptura de pendiente} migran aguas arriba y, como predice el teorema de la onda cinemática para $n
= 1$, se sitúan sobre una misma curva de nivel.
\end{solution}

\begin{solution}{exo:b3:quantitative-geomorphology:5}
$\theta = m/n = 0.5$. $S = (U/K)^{1/2}A^{-1/2} = (10^{-4}/10^{-9})^{1/2}
\times(10^7)^{-1/2} = 316\times3.16\times10^{-4} = 0.10$. Al duplicar $U$
se multiplica $k_s = (U/K)^{1/n}$ por $\sqrt2 = 1.41$; para $n = 1$, por 2:
cuanto mayor es $n$, menos responde la pendiente al levantamiento.
\end{solution}

\begin{solution}{exo:b3:quantitative-geomorphology:6}
Con $x$ aguas arriba desde la desembocadura, $\mathrm{d}x = -\mathrm{d}d$, por lo que $\chi(d)
= \int_d^L (A_0/c\,d'^{h})^{\theta}\,\mathrm{d}d' = A_0^{\theta} c^{-\theta}
\int_d^L d'^{-h\theta}\,\mathrm{d}d' = A_0^{\theta}c^{-\theta}(L^p -
d^p)/p$. Aquí $p = 0.19$, $A_0^{\theta}c^{-\theta} = 10^{2.7}/2^{0.45}$,
y $\chi = \qty{2935}{m}$.
\end{solution}

\begin{solution}{exo:b3:quantitative-geomorphology:7}
$c = 2\times10^{-6}\times(10^8)^{0.5} = \qty{0.02}{m/yr}$, \qty{20}{m/kyr};
A esa velocidad, \qty{5}{km} se recorren en \qty{250}{kyr}. El área, y con ella la
celeridad, disminuye aguas arriba, por lo que el punto de ruptura se frena al ascender.
\end{solution}

\begin{solution}{exo:b3:quantitative-geomorphology:8}
$\bar e = 0.40\times600 = \qty{240}{m}$; las cumbres se elevan en $0.80\times240 =
\qty{192}{m}$; la altitud media desciende $240 - 192 = \qty{48}{m}$.
\end{solution}

\begin{solution}{exo:b3:quantitative-geomorphology:9}
$h = h_0\ln(P_0/E)$: $0.5\ln 6 = \qty{0.90}{m}$ y $0.5\ln1.5 =
\qty{0.20}{m}$. A \qty{0.4}{mm/yr} la erosión supera a la producción máxima,
$P_0$: el suelo se pierde y la ladera queda como roca desnuda, con un
descenso limitado por la meteorización.
\end{solution}

\begin{solution}{exo:b3:quantitative-geomorphology:10}
En \omterm{def:b3:quantitative-geomorphology:steady-profile}{estado estacionario} $\mathrm{d}q/\mathrm{d}x = E$ y $q = 0$ en la divisoria, por
lo que $q = E x$. Entonces $\kappa S = E x\,(1 - S^2/S_c^2)$, una cuadrática cuya
raíz positiva es $S = [-\kappa + \sqrt{\kappa^2 + 4E^2x^2/S_c^2}\,]/(2Ex/S_c^2)$. A
20, 50, 100 y \qty{200}{m}: 0.36, 0.68, 0.89 y 1.03, frente a 0.4, 1.0,
2.0 y 4.0 para la difusión lineal ($S = Ex/\kappa$). La colina tiene una
cresta redondeada, y luego laderas cuyo gradiente tiende a $S_c$: pendientes casi planas,
con un relieve mucho menor que el de la parábola.
\end{solution}

\begin{solution}{exo:b3:quantitative-geomorphology:11}
Según el teorema de la onda cinemática, el cambio llega a un punto de coordenada
$\chi$ tras $\chi/(KA_0^m)$; la red ha respondido por completo cuando ha
alcanzado la última cabecera, de mayor $\chi$. $K A_0^m = 10^{-6}\times
1000 = \qty{e-3}{m/yr}$: el tronco responde en \qty{5.46}{Myr}, la red
en \qty{7.24}{Myr}, del orden de un millón de años, lo que cabía esperar.
\end{solution}

\begin{solution}{exo:b3:quantitative-geomorphology:12}
Aguas arriba de la confluencia, el gradiente estacionario del afluente es $S = (U/2K)A^{-\theta}$,
por lo que $\mathrm{d}z/\mathrm{d}\chi = (U/2K)A_0^{-\theta}$, la mitad que el del tronco.
El punto de ruptura alcanza la confluencia cuando el del tronco está en $\chi_j$; tras un
tiempo adicional $\Delta t$, el del tronco se ha desplazado a $\chi_k$ con $\chi_k -
\chi_j = K A_0^m\,\Delta t$, mientras que el del afluente, con celeridad $2KA^m$,
se ha desplazado a $\chi$ con $\chi - \chi_j = 2KA_0^m\,\Delta t = 2(\chi_k -
\chi_j)$.
\end{solution}

\begin{solution}{pb:b3:quantitative-geomorphology:1}
\textbf{1.} $4.5\times12000^{1.7} = \qty{3.87e7}{m^2}$, \qty{38.7}{km^2};
$4.5\times300^{1.7} = \qty{7.3e4}{m^2}$, \qty{0.073}{km^2}.

\textbf{2.} $\theta \approx -\log(0.007855/0.06641)/\log(37.08/0.388) =
0.928/1.980 = 0.47$.

\textbf{3.} Sí: con una dispersión de unos \qty{10}{\percent} en doce puntos,
$0.466$ no difiere significativamente de 0.5.

\textbf{4.} $0.06641\times\sqrt{388\,000} = \qty{41.4}{m}$;
$0.007855\times\sqrt{3.708\times10^7} = \qty{47.8}{m}$. La media geométrica,
\qty{45.1}{m}, estima $k_s$ para $\theta = 0.5$.

\textbf{5.} Un gradiente es la diferencia de dos altitudes, cada una incierta
de decímetros a metros, dividida por una distancia corta; en tramos suaves
el error es una gran parte de la diferencia. $\chi$ integra en lugar de
derivar.

\textbf{6.} $p = 1 - 1.7\times0.5 = 0.15$; $\chi = (1000/\sqrt{4.5})\,
(12000^{0.15} - 9000^{0.15})/0.15 = \qty{543.1}{m}$.

\textbf{7.} $(193.4 - 170)/543.1 = 0.0431$; $k_{sn} = 0.0431\times A_0^{0.5}
= \qty{43.1}{m}$.

\textbf{8.} La estimación chi: se apoya en altitudes a lo largo de \qty{3}{km} de
río, no en doce gradientes ruidosos.

\textbf{9.} $K = U/k_{sn} = 4.3\times10^{-5}/43.1 = \qty{1.00e-6}{yr^{-1}}$.

\textbf{10.} $10^{-6}\times\sqrt{3.87\times10^7}\times0.00691 =
\qty{4.30e-5}{m/yr}$, \qty{43}{m/Myr}: la desembocadura sigue el ritmo del
levantamiento.

\textbf{11.} En la desembocadura $10^{-6}\times6222 = \qty{6.2}{m/kyr}$; en la
cascada, $d = \qty{7.40}{km}$, $A = \qty{1.70e7}{m^2}$ y
\qty{4.1}{m/kyr}.

\textbf{12.} $170 + 43\times(4.5)^{-0.5}(12000^{0.15} - 300^{0.15})/0.15 =
\qty{405}{m}$: un desnivel de \qty{235}{m}.

\textbf{13.} Para $n = 1$ la ecuación es una onda cinemática: los puntos a ambos
lados del escalón se elevan $U$ por unidad de tiempo, así que el escalón conserva su altura,
es decir, los \qty{25}{m} de la caída.

\textbf{14.} $\chi = (1000/\sqrt{4.5})(12000^{0.15} - 7400^{0.15})/0.15 =
\qty{899.4}{m}$.

\textbf{15.} $t = \chi/(KA_0^{0.5}) = 899.4/10^{-3} = \qty{899}{kyr}$.

\textbf{16.} En $\chi = 899.4 + 300 = \qty{1199}{m}$:
$d^{0.15} = 12000^{0.15} - 1199\times0.15\times\sqrt{4.5}/1000$, $d =
\qty{6.25}{km}$, \qty{5.75}{km} aguas arriba de la desembocadura.

\textbf{17.} En la confluencia ($d = \qty{9.5}{km}$) $\chi = \qty{442.8}{m}$;
la cascada se sitúa donde el valor de $\chi$ del afluente alcanza \qty{899.4}{m}:
$s^{0.15} = 5000^{0.15} - (899.4 - 442.8)\times0.15\times\sqrt{1.5}/1000$,
$s = \qty{4.27}{km}$, \qty{0.73}{km} por encima de la confluencia.

\textbf{18.} Sus confluencias están en $\chi = 1954$ y \qty{999}{m}, más allá de
\qty{899}{m}: el \omterm{def:b3:quantitative-geomorphology:knickpoint}{punto de ruptura de pendiente} aún no las ha alcanzado, así que el cauce principal en ese punto sigue
en su nivel antiguo.

\textbf{19.} $E = 5\times10^{-3}\times0.0086 = \qty{4.3e-5}{m/yr}$,
\qty{43}{m/Myr}: las colinas sobre la cascada se erosionan al ritmo del levantamiento, en
\omterm{def:b3:quantitative-geomorphology:steady-profile}{estado estacionario}.

\textbf{20.} $150^2/(5\times10^{-3}) = \qty{4.5}{Myr}$, cinco veces la edad de
la cascada.

\textbf{21.} El pie de las laderas se rebajó \qty{25}{m} con el cauce: sus
partes bajas tienen pendiente excesiva, quizá deslizándose, mientras que sus crestas siguen
erosionándose al ritmo antiguo. Una onda de empinamiento asciende por cada ladera; tras unos pocos
millones de años recuperan la misma parábola, \qty{25}{m} más abajo.

\textbf{22.} $t = \chi k_{sn}/(U A_0^{0.5})$ varía como $1/U$: de
$899/1.2 = 749$ a $899/0.8 = \qty{1124}{kyr}$.

\textbf{23.} Con $n = 2$ la celeridad crece con el gradiente: la cara empinada
se desplazaría más deprisa que los tramos de su entorno y se mantendría afilada, pero
su velocidad y la edad que se deduce de ella diferirían.

\textbf{24.} El \omterm{def:b3:quantitative-geomorphology:knickpoint}{nivel de base} del río Varden descendió
\textbf{\qty[text-series-to-math=true]{900}{kyr}} hace (\qty{899}{kyr} según el
cálculo).
\end{solution}
